% ============================================================
% Uma trigonometria, cinco percursos
% Texto de trabalho sobre a trigonometria na educação básica
% dos Estados Unidos, Alemanha, França, Rússia e Brasil.
% Compilar: pdflatex -> bibtex -> pdflatex -> pdflatex
% Trechos marcados com % CONFERIR devem ser checados na fonte.
% ============================================================
\documentclass[
    article,
    12pt,
    a4paper,
    sumario=tradicional
]{abntex2}

\usepackage[utf8]{inputenc}
\usepackage[brazil]{babel}
\usepackage{lmodern}
\usepackage[T1]{fontenc}
\usepackage{indentfirst}
\usepackage{microtype}
\usepackage{amsmath,amssymb}
\usepackage{url}
\usepackage{tabularx}
\usepackage{booktabs}
\usepackage{array}

\usepackage[alf,bibjustif]{abntex2cite}

\DeclareMathOperator{\arcsen}{arcsen}

% Ambiente de quadro (ABNT): quadros trazem informação textual
\newcommand{\quadroname}{Quadro}
\newfloat{quadro}{loq}{\quadroname}

\newcolumntype{Y}{>{\raggedright\arraybackslash}X}

\titulo{Uma trigonometria, cinco percursos: o ensino escolar de trigonometria nos Estados Unidos, na Alemanha, na França, na Rússia e no Brasil}
\tituloestrangeiro{One trigonometry, five paths: school trigonometry in the United States, Germany, France, Russia and Brazil}
\autor{[Nome do autor]}
\local{Cuiabá}
\data{2026}

\makeatletter
\hypersetup{
    pdftitle={\@title},
    pdfauthor={\@author},
    pdfsubject={Trigonometria na educação básica: comparação entre países},
    colorlinks=true,
    linkcolor=black,
    citecolor=black,
    urlcolor=blue
}
\makeatother

\setlength{\parindent}{1.3cm}
\setlength{\parskip}{0.2cm}

\begin{document}

\maketitle

\begin{resumoumacoluna}
O texto compara o lugar da trigonometria na educação básica de cinco países
--- Estados Unidos, Alemanha, França, Rússia e Brasil --- a partir de documentos
curriculares oficiais e de livros didáticos de referência. Os cinco
compartilham um mesmo esqueleto, que vai do ângulo à semelhança, da semelhança
às razões no triângulo retângulo, daí ao círculo trigonométrico e às funções
periódicas. As diferenças estão em três pontos: quem decide o currículo (o
país, o estado ou cada rede), em que ano cada etapa aparece e, sobretudo,
quanto espaço recebe a trigonometria que se pode chamar de algébrica, isto é,
identidades, funções inversas e equações. Esse espaço vai do quase
desaparecimento, na Alemanha, a um tratamento sistemático e obrigatório, na
Rússia; França e Estados Unidos o reservam a percursos opcionais, e a França o
faz por meio dos números complexos. O texto registra também curiosidades de
nomenclatura, notação e tradição didática: mnemônicos, o teorema de Al-Kashi, a
notação tg, o grado francês, a origem da palavra seno e os círculos
matemáticos russos. Por fim, discute o que a comparação sugere para o ensino de
trigonometria no Brasil.

\vspace{\onelineskip}
\noindent
\textbf{Palavras-chave}: Trigonometria. Currículo. Educação comparada. Educação básica.
\end{resumoumacoluna}

\renewcommand{\resumoname}{Abstract}
\begin{resumoumacoluna}
\begin{otherlanguage*}{english}
This text compares the place of trigonometry in school mathematics in five
countries --- the United States, Germany, France, Russia and Brazil --- based on
official curriculum documents and reference textbooks. All five share the same
backbone: from angles to similarity, from similarity to the ratios in a right
triangle, and from there to the unit circle and to periodic functions. They
differ in who sets the curriculum, in the school year in which each stage
appears and, above all, in how much room is given to algebraic trigonometry,
that is, identities, inverse functions and equations. That room ranges from
near absence in Germany to a systematic and compulsory treatment in Russia,
while France and the United States leave it to optional tracks, France through
complex numbers. The text also records notational, terminological and
pedagogical curiosities, and closes with a discussion of what the comparison
suggests for the teaching of trigonometry in Brazil.

\vspace{\onelineskip}
\noindent
\textbf{Keywords}: Trigonometry. Curriculum. Comparative education. School mathematics.
\end{otherlanguage*}
\end{resumoumacoluna}

\textual

\section{Introdução}

Poucos assuntos da matemática escolar mudaram tão pouco em conteúdo nos
últimos cem anos quanto a trigonometria. Seno, cosseno, leis dos senos e dos
cossenos, o círculo de raio um e as fórmulas de adição estavam nos livros do
início do século XX e continuam nos de hoje. O que muda de um país para outro
é a ordem, a profundidade e o modo de dar sentido a esses objetos: o seno é
primeiro uma razão entre lados ou uma coordenada no círculo? O radiano é
apresentado como definição ou deduzido de uma proporção? As identidades são
um capítulo para todos ou um apêndice para poucos?

Este texto descreve como Estados Unidos, Alemanha, França e Rússia respondem a
essas perguntas, e usa o Brasil como ponto de referência. A descrição se apoia
nos documentos curriculares oficiais e, onde eles são genéricos, em livros
didáticos de uso amplo. Não se pretende exaustividade. Currículos mudam, e em
países federativos como Estados Unidos e Alemanha a variação entre estados é
grande; o que segue é o percurso típico, não o de cada escola.

Para comparar séries, usa-se a contagem de anos de escolaridade a partir do
primeiro ano obrigatório, por volta dos seis anos de idade. Assim, o \textit{grade
9} americano, a \textit{Klasse 9} alemã, a \textit{troisième} francesa e o 9º ano
brasileiro correspondem, grosso modo, a alunos de 14 a 15 anos. A Rússia tem
onze anos de escola, e a 10ª série russa reúne alunos de cerca de 16 anos, a
idade do 11º ano americano e da 2ª série do ensino médio brasileiro.

\section{Um esqueleto comum}

Antes das diferenças, vale registrar o que se repete. Nos cinco países, a
trigonometria escolar percorre, nesta ordem, as seguintes etapas:

\begin{enumerate}
\item o ângulo como medida e como giro, com o grau e o transferidor;
\item semelhança de triângulos e teorema de Tales, junto com o teorema de
Pitágoras;
\item as razões trigonométricas no triângulo retângulo, justificadas pela
semelhança: a razão entre dois lados depende só do ângulo;
\item o triângulo qualquer, com as leis dos senos e dos cossenos e a fórmula
da área $A = \tfrac12 ab \operatorname{sen}C$;
\item o radiano e o círculo trigonométrico, que estendem seno e cosseno a
qualquer número real;
\item as funções seno e cosseno, seus gráficos e a modelagem de fenômenos
periódicos;
\item a trigonometria algébrica: identidades, fórmulas de adição, funções
inversas, equações e inequações;
\item o Cálculo das funções trigonométricas e o cosseno como ferramenta em
vetores (produto escalar) e, em alguns casos, em números complexos.
\end{enumerate}

As etapas 1 a 3 são universais e concentradas entre o 7º e o 10º ano. As
diferenças aparecem a partir da etapa 4, e ficam marcantes na etapa 7.

\section{Estados Unidos}

Não existe currículo nacional obrigatório nos Estados Unidos. A referência
mais difundida é o \textit{Common Core State Standards for Mathematics}
\cite{ccss2010}, adotado por cerca de quarenta estados a partir de 2010 e, desde
então, revisto ou rebatizado em vários deles. Texas, Virgínia e Flórida, entre
outros, seguem padrões próprios, mas com sequência parecida.

A preparação começa cedo. No 4º ano, o ângulo é apresentado como giro e o grau
como a fração $1/360$ de uma volta completa; no 7º ano entram as relações entre
ângulos complementares, suplementares e opostos pelo vértice; no 8º, a
semelhança por dilatações e o teorema de Pitágoras.

O ensino médio americano costuma seguir a sequência \textit{Algebra I},
\textit{Geometry}, \textit{Algebra II} e \textit{Precalculus}, do 9º ao 12º ano. Em
\textit{Geometry}, o padrão G-SRT.6 é o ponto de partida: por semelhança, as
razões entre os lados de um triângulo retângulo são propriedades do ângulo, o
que leva à definição de seno, cosseno e tangente de um ângulo agudo. Seguem a
relação entre seno e cosseno de ângulos complementares (G-SRT.7) e a resolução
de triângulos retângulos em problemas aplicados (G-SRT.8). As leis dos senos e
dos cossenos e a fórmula da área (G-SRT.9 a 11) aparecem como padrões
avançados, marcados com o sinal (+), destinados a quem pretende cursos
técnicos ou científicos.

Um detalhe da geometria americana merece atenção: o radiano é definido no
capítulo de circunferências (G-C.5), como a constante de proporcionalidade
entre o comprimento do arco e o raio, antes de qualquer menção ao círculo
trigonométrico. Só em \textit{Algebra II} o radiano reaparece como comprimento
de arco no círculo unitário (F-TF.1), e o círculo passa a estender as funções
a todos os números reais (F-TF.2). Nesse mesmo curso entram a modelagem de
fenômenos periódicos com amplitude, frequência e linha média (F-TF.5) e a
identidade $\operatorname{sen}^2\theta + \cos^2\theta = 1$ (F-TF.8).

O restante da trigonometria algébrica fica em \textit{Precalculus}, como padrão
avançado: valores exatos pelos triângulos especiais e simetrias pelo círculo
(F-TF.3 e 4), funções inversas obtidas por restrição de domínio e equações
trigonométricas em contextos de modelagem (F-TF.6 e 7) e a demonstração das
fórmulas de adição (F-TF.9). Desde o ano letivo de 2023--2024 há ainda o
\textit{AP Precalculus} \cite{collegeboard2023}, cuja terceira unidade trata
de funções trigonométricas e polares, incluindo funções senoidais e suas
transformações, inversas, equações, inequações e identidades.

\section{Alemanha}

Na Alemanha a educação é competência dos estados (\textit{Länder}). A
Conferência dos Ministros da Educação, a KMK, fixa padrões nacionais para o
fim do 10º ano \cite{kmk2003} e para o \textit{Abitur}, o exame de conclusão do
ensino médio \cite{kmk2012}, mas cada estado distribui o conteúdo por série no
seu próprio currículo. A descrição a seguir segue o percurso típico do
\textit{Gymnasium}, o ramo que leva à universidade, tomando como exemplo o
currículo da Baviera \cite{isb_lehrplanplus}. Em estados com oito anos de
\textit{Gymnasium} e em estados com nove, os mesmos tópicos podem se deslocar
um ano. % CONFERIR: séries exatas no LehrplanPLUS em vigor.

Os primeiros anos dão grande peso à geometria de construção. O aluno alemão
mede ângulos com o \textit{Geodreieck}, um esquadro com transferidor
embutido, constrói triângulos com régua e compasso e estuda os casos de
congruência. No 9º ano chegam os \textit{Strahlensätze}, versão alemã do
teorema de Tales, e o grupo de Pitágoras: o próprio teorema, o teorema do
cateto e o teorema da altura.

As razões trigonométricas aparecem entre o 9º e o 10º ano, com nomes próprios
para os lados do triângulo retângulo: \textit{Gegenkathete} (cateto oposto),
\textit{Ankathete} (cateto adjacente) e \textit{Hypotenuse}. Em seguida vêm o
círculo unitário (\textit{Einheitskreis}), o radiano (\textit{Bogenmaß}) e a
função seno, estudada sobretudo pelos gráficos de
$f(x) = a\operatorname{sen}\bigl(b(x-c)\bigr) + d$ e por problemas de modelagem,
como a duração do dia ao longo do ano. As equações trigonométricas se limitam,
em geral, a casos simples resolvidos com a calculadora; o arco-seno raramente
é estudado como função. As leis dos senos e dos cossenos, tradicionais no
10º ano, ficaram mais leves ou opcionais em vários estados.

No ensino médio (\textit{Oberstufe}), seno e cosseno entram na Análise, com
derivadas e integrais, e o cosseno reaparece na geometria analítica com
vetores, pela fórmula
\[
\cos\varphi = \frac{\vec a \cdot \vec b}{\lVert \vec a\rVert\,\lVert \vec b\rVert},
\]
que talvez seja o uso mais frequente da trigonometria nas provas do
\textit{Abitur}. As fórmulas de adição (\textit{Additionstheoreme}) e as
identidades em geral perderam espaço na maioria dos currículos, e os números
complexos saíram da escola em quase todos os estados. A calculadora gráfica ou
com sistema de álgebra computacional é comum e, em alguns estados, obrigatória
no exame final.

\section{França}

A França é o caso oposto ao dos Estados Unidos e da Alemanha: tem currículo
nacional único, os \textit{programmes} publicados pelo Ministério da Educação no
\textit{Bulletin officiel}. A escola se divide em primário, \textit{collège}
(da \textit{sixième} à \textit{troisième}, equivalentes ao 6º e ao 9º ano) e
\textit{lycée} (\textit{seconde}, \textit{première} e \textit{terminale}).

No \textit{collège} \cite{men2015}, os ângulos são medidos com transferidor
(\textit{rapporteur}) e a soma dos ângulos do triângulo e os ângulos em retas
paralelas vêm na \textit{cinquième}. O teorema de Pitágoras e sua recíproca
ficam na \textit{quatrième}, série em que, por tradição, o cosseno de um ângulo
agudo entrava sozinho, antes do seno e da tangente. Na \textit{troisième} chegam o
teorema de Tales, a semelhança e as três razões no triângulo retângulo.
% CONFERIR: os programas do collège estão sendo revistos (reforma de 2024--2026).

A reforma do \textit{lycée} de 2019 mudou o lugar da trigonometria. A
\textit{seconde} é comum a todos e usa a trigonometria do triângulo retângulo,
sem funções. A partir da \textit{première}, a matemática completa passou a ser
uma especialidade escolhida pelo aluno \cite{men2019a}. Nela aparece o círculo
trigonométrico, apresentado pelo \textit{enroulement de la droite numérique}: a
reta real é ``enrolada'' sobre o círculo de raio um, e cada número real $x$
corresponde a um ponto, cujas coordenadas definem $\cos x$ e $\operatorname{sen} x$.
Seguem os valores notáveis, a paridade, a periodicidade e os gráficos.

O traço mais característico da trigonometria francesa é o produto escalar. Na
mesma \textit{première} de especialidade, ele é definido de modo a valer
$\vec u \cdot \vec v = \lVert \vec u\rVert\,\lVert \vec v\rVert\cos\theta$, e
dele se deduz a lei dos cossenos, que na França se chama teorema de Al-Kashi.
A lei dos senos tem pouco espaço.

Na \textit{terminale} \cite{men2019b}, a especialidade traz as derivadas de seno
e cosseno, o limite $\operatorname{sen}x/x \to 1$ e as primitivas. A opção
\textit{mathématiques expertes}, voltada a quem segue para as ciências exatas,
trata os números complexos na forma trigonométrica e exponencial, com
$e^{i\theta} = \cos\theta + i\operatorname{sen}\theta$ e as fórmulas de Euler e de
Moivre. É sobretudo por esse caminho que as fórmulas de adição e as identidades
chegam ao aluno francês.

Uma consequência da reforma é discutida na própria França: quem não escolhe a
especialidade termina o \textit{lycée} sem ter visto o radiano nem a função
seno.

\section{Rússia}

\subsection{O currículo}

A Rússia tem currículo nacional: um padrão federal (\textit{FGOS}) e, desde
2023, programas federais únicos por disciplina \cite{fgos}. A partir do 7º
ano, a matemática se divide em duas disciplinas paralelas, Álgebra e
Geometria, cada uma com seus livros. Os mais usados são, em geometria, os de
Atanasyan e colaboradores \cite{atanasyan} e os de Pogorelov e, em álgebra e
análise, os de Kolmogorov e colaboradores \cite{kolmogorov}, de Mordkovich e de
Alimov.

Na geometria do 8º ano, junto com a semelhança e o teorema de Pitágoras, entram
seno, cosseno e tangente do ângulo agudo e seus valores para $30^\circ$,
$45^\circ$ e $60^\circ$. No 9º ano a geometria vai bem mais longe do que nos
outros países nessa idade: seno e cosseno passam a ser definidos para ângulos
de $0^\circ$ a $180^\circ$ por uma semicircunferência unitária, com a identidade
fundamental e as fórmulas de redução para $90^\circ-\alpha$ e $180^\circ-\alpha$;
seguem a fórmula da área, as leis dos senos e dos cossenos, a resolução de
triângulos e o produto escalar de vetores. Tudo isso é conteúdo de todos os
alunos e cai no exame de fim do 9º ano.

O núcleo da trigonometria russa está na disciplina Álgebra e Princípios de
Análise do 10º ano. O radiano e a ``circunferência numérica''
(\textit{chislovaya okruzhnost'}), em espírito parecido com o
\textit{enroulement} francês, abrem o capítulo; depois vêm seno, cosseno,
tangente e cotangente para todo número real; as fórmulas de adição, de arco
duplo e de arco metade, as fórmulas de redução e a transformação de somas em
produtos; os gráficos e as propriedades das funções; arco-seno, arco-cosseno e
arco-tangente como funções, e não apenas como teclas da calculadora; e, por
fim, equações e inequações trigonométricas, com fórmulas gerais como
\[
\operatorname{sen}x = a \iff x = (-1)^k\arcsen a + \pi k,\quad k\in\mathbb{Z}\qquad (|a|\le 1),
\]
e métodos de substituição, fatoração e equações homogêneas. No 11º ano vêm as
derivadas e as integrais. O exame de saída de nível avançado tem uma questão
típica de equação trigonométrica com seleção das raízes num intervalo, e essa
questão orienta boa parte do ensino nos dois últimos anos.

\subsection{Os círculos matemáticos}

A fama russa em matemática escolar vem menos do currículo do que de uma
tradição paralela a ele: os círculos matemáticos (\textit{matematicheskie
kruzhki}). Eles funcionam fora da grade, em universidades, nos antigos
Palácios de Pioneiros, em escolas especializadas em física e matemática ---
como a escola-internato criada por Kolmogorov junto à Universidade de Moscou
em 1963 --- e no ensino por correspondência fundado por Gelfand. Sustentam-se
nas olimpíadas e na revista \textit{Kvant}, e o relato mais conhecido dessa
experiência, escrito a partir dos círculos de Leningrado, está em
\citeonline{fomin1996}.

Nos círculos, a trigonometria não é conteúdo a cumprir, e sim ferramenta para
resolver problemas de geometria. O chamado método trigonométrico em
planimetria usa a lei dos senos, a forma trigonométrica do teorema de Ceva e
identidades do triângulo, como $\cos A + \cos B + \cos C = 1 + r/R$, em que $r$ e
$R$ são os raios das circunferências inscrita e circunscrita. O livro de
\citeonline{gelfand2001}, nascido do ensino por correspondência, mostra o
espírito: a trigonometria é tratada como algo a ser entendido a partir de
poucas ideias, e não como uma coleção de fórmulas.

\section{O Brasil como ponto de referência}

O Brasil se parece com a França e com a Rússia na organização: há uma base
nacional, a BNCC \cite{bncc2018}, que as redes estaduais e municipais
detalham. No Ensino Fundamental, a BNCC chega à semelhança e ao teorema de
Pitágoras no 9º ano (habilidades EF09MA13 e EF09MA14). No Ensino Médio, que a
BNCC organiza por competências e não por série, a trigonometria aparece de
forma explícita em habilidades como EM13MAT306, sobre fenômenos periódicos e
as funções seno e cosseno, e EM13MAT308, que inclui as leis dos senos e dos
cossenos. % CONFERIR: códigos e redação das habilidades na versão oficial.
A série em que as razões trigonométricas, o círculo e as funções aparecem
depende do currículo de cada rede e do livro adotado.

\section{Quadro comparativo}

O Quadro~\ref{quadro:comparativo} resume a comparação. As séries indicadas são
as típicas e podem variar dentro de cada país.

\begin{quadro}[htb]
\footnotesize
\caption{A trigonometria escolar em cinco países: síntese}
\label{quadro:comparativo}
\begin{tabularx}{\textwidth}{@{}>{\raggedright\arraybackslash}p{2.4cm}YYYYY@{}}
\toprule
 & \textbf{EUA} & \textbf{Alemanha} & \textbf{França} & \textbf{Rússia} & \textbf{Brasil}\\
\midrule
Quem decide & estados (Common Core como referência) & estados (padrões da KMK) & país & país & país (BNCC), detalhada pelas redes\\
Razões no triângulo retângulo & 10º ano & 9º--10º ano & 9º ano & 8º ano & em geral 9º ano\\
Leis dos senos e cossenos & avançado (+) & 10º ano, mais leve & Al-Kashi pelo produto escalar & 9º ano, para todos & ensino médio\\
Radiano e círculo & radiano pela proporção do arco; círculo no 11º & 10º ano & especialidade da \textit{première} & 10º ano & ensino médio\\
Identidades e adição & avançado (Precalculus) & quase ausentes & via complexos, opção \textit{expertes} & 10º ano, para todos & varia muito\\
Inversas e equações & avançado & só casos simples & pouco espaço & 10º ano, em profundidade & varia muito\\
Cálculo & AP Calculus & \textit{Oberstufe} & \textit{terminale} & 11º ano & em geral ausente\\
Notação da tangente & tan & tan & tan & tg & tg\\
\bottomrule
\end{tabularx}
\fonte{Elaboração própria a partir das fontes citadas no texto.}
\end{quadro}

\section{Curiosidades}

\subsection{Mnemônicos}
 Nos Estados Unidos, as três razões se decoram com
``SOH-CAH-TOA'' (\textit{sine} = \textit{opposite}/\textit{hypotenuse}, e assim por
diante). Na França, a mesma sigla vira ``CAH-SOH-TOA'', que se lê como
\textit{casse-toi}, um ``some daqui'' que nenhum aluno esquece. Na Alemanha, um
mnemônico conhecido é ``GAGA Hühnerhof AG'': as letras G, A, G, A são os
numeradores e as iniciais seguintes, H, H, A, G, os denominadores, o que dá
$\operatorname{sen} = G/H$, $\cos = A/H$, $\tan = G/A$ e $\cot = A/G$, com $G$ de
\textit{Gegenkathete}, $A$ de \textit{Ankathete} e $H$ de \textit{Hypotenuse}.

\subsection{Al-Kashi}
 O nome francês da lei dos cossenos homenageia o
astrônomo e matemático persa Jamshid al-Kashi, do século XV, que a usou de
forma sistemática em cálculos astronômicos.

\subsection{tg ou tan?}
 A notação também separa tradições. Estados Unidos,
Alemanha e França escrevem tan e cot. Rússia e Brasil escrevem tg, e para a cotangente o Brasil
usa cotg, onde os russos escrevem ctg; arco-seno se escreve
arcsen em português e arcsin em russo, inglês, alemão e francês.

\subsection{O grado}
 As calculadoras científicas têm três modos de ângulo:
DEG, RAD e GRAD. O terceiro é o grado (ou \textit{gon}), em que uma volta vale
400 unidades e o ângulo reto vale 100. Ele nasceu na França, no espírito
decimal do sistema métrico, não se firmou na escola e sobrevive na topografia.

\subsection{Por que ``seno''?}
 A palavra tem uma das etimologias mais tortas da
matemática \cite{maor1998,vanbrummelen2009}. Os astrônomos indianos chamavam a
meia-corda de \textit{jya}; os árabes transliteraram o termo como \textit{jiba},
escrito só com consoantes. Tradutores latinos leram a palavra como
\textit{jaib}, que significa dobra ou enseada, e a traduziram por
\textit{sinus}, de onde vêm o seno, o \textit{sine} e o \textit{Sinus}.

\subsection{Definir ou deduzir o radiano}
 Nos Estados Unidos, o radiano nasce
de uma proporção (o arco é proporcional ao raio); na França e na Rússia, de um
enrolamento da reta sobre o círculo; na Alemanha, como uma nova unidade de
medida associada ao círculo unitário. São três modos de dar sentido ao mesmo
número, e cada um sugere uma atividade diferente em sala de aula.

\section{Considerações finais}

A comparação mostra um mesmo esqueleto e uma divergência concentrada no que se
faz depois da função seno. A Alemanha praticamente abandonou a trigonometria
algébrica e investiu em modelagem, Cálculo e vetores. A França a manteve, mas
em ramo opcional e apoiada nos números complexos. Os Estados Unidos a
reservaram a padrões avançados. A Rússia a trata para todos, em profundidade,
no 10º ano, e mantém nos círculos uma trigonometria que nasce de problemas.

Para o Brasil, duas observações parecem úteis. A primeira é que a ordem
``medir, prever e só então animar'', que aparece com clareza nos documentos
americanos e alemães, tem respaldo internacional: o círculo trigonométrico
funciona melhor como ponte para razões já compreendidas do que como ponto de
partida. A segunda é que a ausência de identidades e equações como eixo de
pesquisa e de ensino não é uma fatalidade. A experiência russa mostra que esses
tópicos podem ocupar o centro do currículo, e a francesa mostra que podem ser
apresentados como consequência de uma estrutura mais rica, a dos números
complexos ou do produto escalar.

\section*{Nota sobre as fontes}
\addcontentsline{toc}{section}{Nota sobre as fontes}

A descrição dos currículos é um retrato do percurso típico em cada país e
deve ser conferida na versão em vigor dos documentos citados antes de ser usada
como dado em pesquisa, em especial a distribuição por série na Alemanha e na
Rússia, os códigos da BNCC e os programas franceses do \textit{collège}, que estão
em revisão. Os códigos do \textit{Common Core} seguem a numeração oficial de 2010.

\postextual

\bibliography{referencias}

\end{document}
